\begin{frame}{자주 묻는 질문 — GAE와 PPO의 하이퍼파라미터}
  \small
  \begin{itemize}
    \item \kb{Q. GAE의 $\lambda$ 와 클리핑의 $\epsilon$ 은
      독립적?} 역할은 다르지만 \textbf{함께} 학습 안정성에 영향 —
      $\lambda$ 는 ``어드밴티지 추정 \textit{자체}의
      편향-분산'', $\epsilon$ 은 ``그 추정을 바탕으로 정책을
      얼마나 크게 바꿀 것인가''. 어드밴티지 추정이 노이즈가
      크면($\lambda$ 가 1에 가까우면) 클리핑이 그 노이즈로부터
      정책을 \textbf{보호}하는 역할이 더 중요 — 실전에서는 둘
      함께 튜닝.
    \item \kb{Q. RLHF에서 PPO는 무엇을 ``행동''으로 보나?}
      언어모델을 정책으로 볼 때 \kb{상태} = 지금까지 생성된
      토큰들, \kb{행동} = 다음에 고를 토큰, \kb{보상} =
      보상모델의 점수. 토큰 하나하나를 고르는 것이 이산 행동
      선택이므로 이번 장의 \textbf{이산 행동공간 PPO가 사실상
      그대로 적용} — 로봇(연속)과 언어모델(이산)이라는 완전히
      다른 응용에 같은 알고리즘이 쓰이는 것 자체가 ``행동공간
      형태에 크게 의존하지 않는'' 범용성의 증거.
    \item \kb{Q. GAE의 $\lambda$ 와 7절 적격흔적 TD($\lambda$)의
      $\lambda$ 는 같은 것인가?} ``같은 아이디어''지만
      \textbf{적용 대상이 다름} — 적격흔적의 $\lambda$ 는
      \textit{가치함수 $V$}의 갱신에 ``어떤 $n$-step TD를
      가중합할지'', GAE의 $\lambda$ 는 \textit{어드밴티지
      $A_t$}를 ``어떤 $n$-step 어드밴티지를 가중합할지'' 결정.
      수학적 구조(지수 가중 $\lambda^k$ 의 합)는 같음 —
      적격흔적은 critic 학습에, GAE는 actor의 어드밴티지
      계산에.
    \item \kb{Q. 엔트로피 $c_e$ 를 너무 크게?} ``무작위에
      가까운 정책''이 장려되어, 좋은 행동을 찾았을 때도 확률을
      높이지 못해(엔트로피 최대화가 ``균등한 확률''을 유지하려
      함) 성능 오히려 나빠짐 — 실전 $c_e=0.001\sim0.01$ 로
      작게, 끝단계에서 0으로 줄여(annealing) ``탐험 종료''.
  \end{itemize}
\end{frame}
